\documentclass[11pt]{article}
\usepackage[letterpaper,margin=0.85in]{geometry}\usepackage[T1]{fontenc}\usepackage[utf8]{inputenc}\usepackage{amsmath,amssymb,enumitem,booktabs,array,xcolor,fancyhdr,tabularx}
\setlist{itemsep=2pt,topsep=3pt}\pagestyle{fancy}\fancyhf{}\renewcommand{\headrulewidth}{0.3pt}
\fancyhead[L]{\small 4.4 Part 1 $\cdot$ Related Rates: Set It Up}\fancyhead[R]{\small Teacher key}\fancyfoot[C]{\small\thepage}
\definecolor{ans}{RGB}{28,125,74}\newcommand{\A}[1]{\textcolor{ans}{\textbf{#1}}}
\begin{document}

\section*{4.4 Part 1 $\cdot$ Related Rates: Set It Up --- teacher key}
\textbf{Question:} When two things change together, how do I name them and link them?\par
\subsection*{Self-paced lesson}No clock, no trimming, no skipping, no end mark, no stamp. Every tab is required in order; progress saves on the device so a student can stop and finish the next day.\par
\subsection*{Warm-up review (3 easy items)}\begin{itemize}
\item The derivative of $x^3$ is... \A{$3x^2$.} Power rule: bring down the 3 and lower the power by one.
\item If $y$ depends on $x$, then $\dfrac{d}{dx}(y^2)$ is... \A{$2y\dfrac{dy}{dx}$.} Chain rule: differentiate $y^2$ to get $2y$, then multiply by $\dfrac{dy}{dx}$.
\item A particle has velocity $v(t)=-3$ at one moment. It is moving... \A{in the negative direction.} A negative velocity means the position is getting smaller.
\end{itemize}
\subsection*{Independent lesson}No huddle, no teacher talk, no codes. Students read a short picture tutorial (Concept), watch worked problems move by move (I do), then solve generated problems (You do). Add \texttt{\#key} to the lesson address to open every tab with answers marked.\par
\subsection*{Targets}\begin{itemize}\item I can name the quantities that change together in a story.
\item I can list the rates I know and the rate I want, with signs.
\item I can pick the equation that links the quantities.\end{itemize}
\subsection*{Predictions}\begin{itemize}
\item A balloon is blown up so its radius grows by the same amount every second. Its volume grows... \A{by more and more each second.} Volume depends on $r^3$, so the bigger the balloon, the faster its volume grows. The tutorial shows why.
\item A ladder's foot is pulled away from a wall at a steady speed. The top of the ladder slides down... \A{slowly at first, then faster and faster.} The top barely moves at first and drops quickly at the end. You will see the reason in 4.5.
\end{itemize}
\section*{Tutorial script, with every question and answer}
\subsection*{Concept tutorial (8 pictures)}
\textit{I can say which quantities change, what a rate is, and which equation links them.}\par
\begin{enumerate}
\item \textbf{Two things change at once.} As air goes in, the radius $r$ and the volume $V$ both change. They change together. \textit{(Picture: balloon. A balloon from the side. The line from the center to the edge is the radius $r$. Slide $r$ and read $V$ under the picture.)}\\ \textit{Question:} Which quantities change as the balloon fills? \A{Answer: Both $r$ and $V$.} Both change. When one changes, the other is forced to change with it.
\item \textbf{Not at the same speed.} Each time $r$ goes up by 1 cm, $V$ grows by a bigger amount than the time before. \textit{(Picture: balloon. The table lists some radii and the volume for each. The gold row follows the slider. Compare how much $V$ grows each time $r$ goes up by 1.)}\\ \textit{Question:} If the radius grows at a steady pace, the volume grows... \A{Answer: faster and faster.} V goes 4.2, 33.5, 113.1, 268.1. The jumps are 29, 80, 155: bigger each time.
\item \textbf{A ladder ties two lengths together.} The foot slides out and the top slides down. The ladder never gets longer, so $x$ and $y$ are tied together. \textit{(Picture: ladder. The brown line is a ladder 5 ft long. $x$ is how far the foot is from the wall. $y$ is how high the top is. Slide $x$ and watch $y$.)}\\ \textit{Question:} When the foot is $x=3$ ft from the wall, how high is the top? \A{Answer: 4 ft.} $3^2+y^2=25$, so $y^2=16$ and $y=4$ ft.
\item \textbf{A rate says how fast.} The symbol $\dfrac{dx}{dt}$ means how fast $x$ changes each second. Its units are feet per second. \textit{(Picture: ladder. The red arrows show which way each length is moving. $\dfrac{dx}{dt}$ is how fast $x$ changes as time $t$ passes.)}\\ \textit{Question:} $\dfrac{dy}{dt}=-2$ ft/s. What does that say? \A{Answer: The top is sliding down 2 feet each second..} The 2 is how fast. The minus sign says $y$ is getting smaller.
\item \textbf{Growing is positive. Shrinking is negative..} If a quantity is growing, its rate is positive. If it is shrinking, its rate is negative. \textit{(Picture: cylinder. A tank with water draining out. The red label says the volume is getting smaller. The slider sets the depth $h$.)}\\ \textit{Question:} Water drains out of the tank. The rate $\dfrac{dV}{dt}$ is... \A{Answer: negative, because $V$ is shrinking.} Water leaving means $V$ is getting smaller, so $\dfrac{dV}{dt}<0$.
\item \textbf{One equation links them.} A link equation joins the changing quantities, like $A=\pi r^2$. A number like $\pi$ never changes, so it stays. \textit{(Picture: formulas. Six shapes and the equation for each. The blue card is the circle: area $A$ and radius $r$.)}\\ \textit{Question:} A circular puddle has area $A$ and radius $r$. Which equation links them? \A{Answer: $A=\pi r^2$.} The area of a circle is $\pi r^2$.
\item \textbf{Given and wanted.} "The radius grows 3 ft per minute" gives $\dfrac{dr}{dt}=3$. "How fast is the area growing?" asks for $\dfrac{dA}{dt}$. \textit{(Picture: balloon. A puddle seen from above. The red arrows show its radius growing.)}\\ \textit{Question:} The radius grows 3 ft/min. We are asked how fast the area grows. Which rate do we WANT? \A{Answer: $\dfrac{dA}{dt}$.} We are given $\dfrac{dr}{dt}=3$. The thing we want is $\dfrac{dA}{dt}$.
\item \textbf{Which numbers go in the equation?.} Only numbers that never change may go into the link equation. Changing quantities stay as letters. \textit{(Picture: ladder. A 13 ft ladder. The foot $x$ and the top $y$ change as it slides. The 13 ft never changes.)}\\ \textit{Question:} The link is $x^2+y^2=13^2$. Why is 13 allowed, but "$x=5$" is not? \A{Answer: 13 never changes. $x$ changes as the ladder slides..} The ladder length stays 13 ft. The foot distance $x=5$ is true only at one moment.
\end{enumerate}
\subsection*{I do: the five moves}\begin{enumerate}\item Draw and name
\item List what you know and want
\item Link with one equation
\item Differentiate in t
\item Plug in last\end{enumerate}
\subsection*{I do (8 steps)}
\textit{I can watch the first three moves done on two problems and spot the common slip.}\par
\begin{enumerate}
\item \textbf{Name what changes.} \textit{Move 1 $\cdot$ Draw and name.} Name every quantity that changes: the radius $r$ (cm), the area $A$ (cm\textsuperscript{2}) and time $t$ (s). \textit{(Picture: balloon. The puddle from above. The red arrows show the radius growing.)}\\ Changing: $r$ (cm), $A$ (cm\textsuperscript{2}), $t$ (s)\\ \textit{Question:} Which of these changes while the puddle spreads? \A{Answer: $r$ and $A$.} $r$ and $A$ both change. $\pi$ never changes, and 10 cm is the radius at just one moment.
\item \textbf{Write what we know and what we want.} \textit{Move 2 $\cdot$ List what you know and want.} Write each rate with its symbol. "Grows 2 cm every second" is a rate. \textit{(Picture: balloon. Same puddle. Growing means a positive rate.)}\\ Given: $\dfrac{dr}{dt}=2$ cm/s\\ Want: $\dfrac{dA}{dt}$ when $r=10$ cm\\ \textit{Question:} How do we write "the radius grows 2 cm every second"? \A{Answer: $\dfrac{dr}{dt}=2$.} A rate of change in $r$, and growing makes it positive.
\item \textbf{The moment waits.} \textit{Move 2 $\cdot$ List what you know and want.} The 10 cm is true only at this one moment. We keep it on the side until the last move. \textit{(Picture: balloon. The puddle at the moment we are asked about: $r=10$ cm. This is one frozen moment.)}\\ Save for the end: $r=10$\\ \textit{Question:} Should we put $r=10$ into the equation now? \A{Answer: No. The radius keeps changing, so 10 is true for one moment only..} Numbers for one moment wait. Only numbers that never change may go in early.
\item \textbf{Link with one equation.} \textit{Move 3 $\cdot$ Link with one equation.} The equation that joins $A$ and $r$ is the area of a circle, and we keep $r$ as a letter. \textit{(Picture: formulas. The circle card is highlighted: area and radius.)}\\ Link: $A=\pi r^2$\\ \textit{Question:} Which link equation do we write now? \A{Answer: $A=\pi r^2$.} Keep $r$ as a letter. We will use 10 in the last move.
\item \textbf{Problem 2: a sliding ladder.} \textit{Move 1 $\cdot$ Draw and name.} \textbf{Problem:} A 13 ft ladder slides away from a wall. The foot moves out at 2 ft/s. How fast is the top falling when the foot is 5 ft from the wall? Name what changes: the foot $x$ and the top $y$. The ladder $L=13$ ft does not change. \textit{(Picture: ladder. A 13 ft ladder against a wall. $x$ is the foot, $y$ is the top.)}\\ Changing: $x$ (ft), $y$ (ft)\\ Never changes: $L=13$ ft\\ \textit{Question:} Which quantities change as the ladder slides? \A{Answer: $x$ and $y$.} The foot and the top both move. The ladder itself keeps its length.
\item \textbf{Write what we know and what we want.} \textit{Move 2 $\cdot$ List what you know and want.} The foot moves out at 2 ft/s. We want the rate of the top, which is falling. \textit{(Picture: ladder. The red arrows show the foot moving out and the top moving down.)}\\ Given: $\dfrac{dx}{dt}=2$ ft/s\\ Want: $\dfrac{dy}{dt}$ when $x=5$ ft\\ \textit{Question:} The top is FALLING. What sign should $\dfrac{dy}{dt}$ have? \A{Answer: negative.} $y$ is getting smaller, so its rate is negative.
\item \textbf{Link with one equation.} \textit{Move 3 $\cdot$ Link with one equation.} The ladder, the wall and the floor make a right triangle. The 13 may stay, because it never changes. \textit{(Picture: ladder. The ladder is the long side of a right triangle.)}\\ Link: $x^2+y^2=13^2$\\ \textit{Question:} Which number may go into the link equation now? \A{Answer: 13, the ladder length.} 13 never changes. The 5 and 12 are for one moment, and 2 is a rate, not part of the link.
\item \textbf{The common slip.} Slip: writing $5^2+y^2=13^2$. A fixed 5 cannot change, so $\dfrac{dx}{dt}$ has nothing to belong to. \textit{(Picture: ladder. The foot is labeled with the fixed number 5. A number that is fixed cannot have a rate.)}\\ Wrong: $5^2+y^2=13^2$\\ Right: $x^2+y^2=13^2$\\ \textit{Common slip.}\\ \textit{Question:} Which equation should we carry into the next part? \A{Answer: $x^2+y^2=13^2$.} Keep $x$ and $y$ as letters. Next time we differentiate, then use 5 and 12.
\end{enumerate}
\subsection*{Big ideas}\begin{itemize}\item Quantities in a story change with time, so each one has a rate, like $\dfrac{dr}{dt}$.
\item Growing means a positive rate. Shrinking means a negative rate.
\item One link equation joins the changing quantities. Only numbers that never change go in it.
\item Numbers true at one moment (like $r=10$) wait until the last move.\end{itemize}
\section*{Every level, with answers}
\subsection*{Pick the link equation}
\noindent\textbf{Level 1.} A spherical balloon has radius $r$ and volume $V$. Which equation links the changing quantities?\\ \A{(A) $V=\tfrac43\pi r^3$} \quad (B) $V=4\pi r^2$ \quad (C) $V=\pi r^2$ \quad (D) $V=\tfrac13\pi r^3$\\ \A{Answer: A} \quad The volume of a sphere is $\tfrac43\pi r^3$. The $4\pi r^2$ is its surface area. So the link is $V=\tfrac43\pi r^3$.\par\medskip
\noindent\textbf{Level 2.} A square window has side $s$ and area $A$. At one moment $s=3$ in. Which equation links the changing quantities?\\ (A) $A=2s$ \quad (B) $A=4s$ \quad (C) $A=3^2$ \quad \A{(D) $A=s^2$}\\ \A{Answer: D} \quad A square's area is side times side. The $4s$ is the distance around it. So the link is $A=s^2$, and the 3 waits for the last move.\par\medskip
\noindent\textbf{Level 3.} A plane flies level 3 miles up. It is $x$ miles along the ground from a station and $s$ miles from it in a straight line. At one moment $x=3$ mi. Which equation links the changing quantities?\\ (A) $x^2+s^2=3^2$ \quad (B) $3^2+3^2=s^2$ \quad \A{(C) $x^2+3^2=s^2$} \quad (D) $x+3=s$\\ \A{Answer: C} \quad The height 3 never changes, so it may stay in the equation as a number. So the link is $x^2+3^2=s^2$, and the 3 waits for the last move.\par\medskip
\subsection*{Given and wanted}
\noindent\textbf{Level 1.} A balloon is being filled so that its volume grows 3 cm\textsuperscript{3} every second. How fast is its radius growing when the radius is 10 cm? Which line sets it up?\\ (A) Given $\dfrac{dV}{dt}=10$. Want $\dfrac{dr}{dt}$. \quad \A{(B) Given $\dfrac{dV}{dt}=3$. Want $\dfrac{dr}{dt}$.} \quad (C) Given $\dfrac{dr}{dt}=3$. Want $\dfrac{dV}{dt}$. \quad (D) Given $\dfrac{dV}{dt}=-3$. Want $\dfrac{dr}{dt}$.\\ \A{Answer: B} \quad The rate we are told is $\dfrac{dV}{dt}=3$ (growing, so positive). The rate we are asked for is $\dfrac{dr}{dt}$. The 10 is only for the last move.\par\medskip
\noindent\textbf{Level 2.} A ladder's foot slides away from a wall at 6 ft/s. How fast is the top falling when the foot is 11 ft from the wall? Which line sets it up?\\ (A) Given $\dfrac{dx}{dt}=11$. Want $\dfrac{dy}{dt}$. \quad (B) Given $\dfrac{dx}{dt}=-6$. Want $\dfrac{dy}{dt}$. \quad (C) Given $\dfrac{dy}{dt}=6$. Want $\dfrac{dx}{dt}$. \quad \A{(D) Given $\dfrac{dx}{dt}=6$. Want $\dfrac{dy}{dt}$.}\\ \A{Answer: D} \quad The rate we are told is $\dfrac{dx}{dt}=6$ (growing, so positive). The rate we are asked for is $\dfrac{dy}{dt}$. The 11 is only for the last move.\par\medskip
\noindent\textbf{Level 3.} The top of a ladder slides down a wall at 2 ft/s. How fast is the foot moving out when the foot is 9 ft from the wall? Which line sets it up?\\ (A) Given $\dfrac{dy}{dt}=2$. Want $\dfrac{dx}{dt}$. \quad (B) Given $\dfrac{dx}{dt}=-2$. Want $\dfrac{dy}{dt}$. \quad (C) Given $\dfrac{dy}{dt}=9$. Want $\dfrac{dx}{dt}$. \quad \A{(D) Given $\dfrac{dy}{dt}=-2$. Want $\dfrac{dx}{dt}$.}\\ \A{Answer: D} \quad The rate we are told is $\dfrac{dy}{dt}=-2$ (shrinking, so negative). The rate we are asked for is $\dfrac{dx}{dt}$. The 9 is only for the last move.\par\medskip
\subsection*{Sign of a rate}
\noindent\textbf{Level 1.} A snowman's height is shrinking 4 cm each hour. Which line writes its rate?\\ \A{(A) $\dfrac{dh}{dt}=-4$} \quad (B) $\dfrac{dh}{dt}=4$ \quad (C) $\dfrac{dt}{dh}=-4$ \quad (D) $h=-4$\\ \A{Answer: A} \quad The quantity is $h$ and it is shrinking, so the rate is negative: $\dfrac{dh}{dt}=-4$.\par\medskip
\noindent\textbf{Level 2.} A tank is being drained at 8 gallons per minute. Which line writes its rate?\\ (A) $\dfrac{dV}{dt}=8$ \quad (B) $V=-8$ \quad \A{(C) $\dfrac{dV}{dt}=-8$} \quad (D) $\dfrac{dt}{dV}=-8$\\ \A{Answer: C} \quad The quantity is $V$ and it is shrinking, so the rate is negative: $\dfrac{dV}{dt}=-8$.\par\medskip
\noindent\textbf{Level 3.} A shadow is getting longer by 2 ft each second. Which line writes its rate?\\ (A) $\dfrac{dt}{ds}=2$ \quad (B) $\dfrac{ds}{dt}=-2$ \quad \A{(C) $\dfrac{ds}{dt}=2$} \quad (D) $s=2$\\ \A{Answer: C} \quad The quantity is $s$ and it is growing, so the rate is positive: $\dfrac{ds}{dt}=2$.\par\medskip
\subsection*{What goes in the equation now?}
\noindent\textbf{Level 1.} A puddle's radius is growing. Right now $r=6$ cm. The link is $A=\pi r^2$. What may go into this equation BEFORE we differentiate?\\ \A{(A) Only $\pi$ (it never changes). Not $r=6$.} \quad (B) Every number in the story, including the one for this moment. \quad (C) The rates too, like $\dfrac{dr}{dt}$. \quad (D) Nothing at all, not even $\pi$ or the 13.\\ \A{Answer: A} \quad Changing quantities stay as letters until the last move. Only $\pi$ (it never changes). Not $r=6$.\par\medskip
\noindent\textbf{Level 2.} A can has a fixed radius 4 cm and water of changing depth. Right now $h=4$ cm. The link is $V=\pi r^2h$. What may go into this equation BEFORE we differentiate?\\ (A) Every number in the story, including the one for this moment. \quad (B) Nothing at all, not even $\pi$ or the 13. \quad \A{(C) The fixed radius and $\pi$. Not $h=4$.} \quad (D) The rates too, like $\dfrac{dr}{dt}$.\\ \A{Answer: C} \quad Changing quantities stay as letters until the last move. The fixed radius and $\pi$. Not $h=4$.\par\medskip
\noindent\textbf{Level 3.} A 13 ft ladder is sliding. Right now $x=5$ ft and $y=12$ ft. The link is $x^2+y^2=13^2$. What may go into this equation BEFORE we differentiate?\\ (A) The rates too, like $\dfrac{dr}{dt}$. \quad (B) Every number in the story, including the one for this moment. \quad \A{(C) Only the 13. Not $x=5$ and not $y=12$.} \quad (D) Nothing at all, not even $\pi$ or the 13.\\ \A{Answer: C} \quad Changing quantities stay as letters until the last move. Only the 13. Not $x=5$ and not $y=12$.\par\medskip
\section*{Common misconceptions}\begin{tabularx}{\textwidth}{@{}X X@{}}\toprule Students say & Say instead\\\midrule
Put r = 10 into A = $\pi$r\textsuperscript{2} first. & Keep r as a letter. 10 is true only at one moment, so it waits for the last move.\\
The rate is the same as the quantity (r = 2 means dr/dt = 2). & dr/dt says how fast r changes. r is how big it is.\\
A shrinking quantity has a positive rate. & Shrinking means the rate is negative.\\
Using the area of a circle for the circumference (or the reverse). & Area is $\pi$r\textsuperscript{2}. The distance around is 2$\pi$r.\\
Putting a changing number (like the foot distance) into the link equation. & Only numbers that never change (like the ladder length) may go in.\\
\bottomrule\end{tabularx}
\end{document}